\chapter{Conclusioni}

Il progetto ha dimostrato l'efficacia dell'Hybrid Genetic Algorithm per la
risoluzione del Capacitated Vehicle Routing Problem. L'approccio ibrido, che
combina la potenza esplorativa degli algoritmi genetici con l'efficacia della
ricerca locale, produce soluzioni di alta qualit\`a sulle istanze benchmark
CVRPLIB.

\section{Scelte Progettuali e Originalit\`a}

\subsection{Scelte Chiave}

\begin{enumerate}[label=\textbf{\arabic*.}]
    \item \textbf{Split Algorithm} \cite{prins2004simple}: L'uso dello split di Prins come decodifica
    \`e fondamentale. Invece di codificare esplicitamente le route nel cromosoma,
    si lascia che la DP trovi la partizione ottimale, riducendo lo spazio di
    ricerca e garantendo l'ottimalit\`a della decodifica data la permutazione.

    \item \textbf{Ricerca Locale Granulare}: L'integrazione della GLS con
    $\gamma = 25$ permette di esplorare un vicinato inter-route molto pi\`u
    ampio rispetto alle configurazioni standard ($\gamma = 15$), senza il
    costo computazionale proibitivo della ricerca esaustiva $O(N^2)$.

    \item \textbf{Ottimizzazione Iperparametrica}: L'uso di Optuna per il
    tuning automatico dei parametri ha permesso di scoprire una configurazione
    non ovvia ($\mu = 81$, $p_c = 0.675$, $p_m = 0.236$, $p_{ls} = 0.259$)
    che supera le configurazioni progettate manualmente.

    \item \textbf{Accelerazione Numba}: L'uso del compilatore JIT Numba per i
    calcoli critici (matrici di distanza, Split DP, operatori di ricerca locale)
    ha permesso di ottenere performance computazionali elevate in Python puro,
    senza ricorrere a implementazioni in C/C++.

    \item \textbf{Architettura client-server con WebSocket}: L'algoritmo \`e
    integrato in un'architettura moderna con backend Python (FastAPI) e frontend
    React, permettendo la visualizzazione in tempo reale dell'evoluzione
    dell'algoritmo tramite WebSocket.
\end{enumerate}

\subsection{Difficolt\`a Affrontate}

\begin{itemize}
    \item \textbf{Vincoli di capacit\`a}: Gestire i vincoli di capacit\`a durante
    il crossover e la mutazione \`e complesso perch\'e questi operatori lavorano
    sulla permutazione, non sulle route. Lo split algorithm risolve elegantemente
    questo problema.

    \item \textbf{Bilanciamento esplorazione-sfruttamento}: Trovare il giusto
    tasso di ricerca locale \`e stato cruciale: troppo alto causa convergenza
    prematura, troppo basso rallenta il miglioramento. Il tuning con Optuna ha
    permesso di individuare il valore ottimale $p_{ls} = 0.259$.

    \item \textbf{Scalabilit\`a}: Su istanze con 101 nodi, lo split $O(n^2)$ e
    la ricerca locale diventano computazionalmente significativi, richiedendo
    attenzione all'efficienza implementativa e all'uso di Numba JIT.
\end{itemize}

\section{Risultati Raggiunti}

I punti di forza dell'implementazione includono:
\begin{itemize}
    \item L'uso dello split algorithm per una decodifica ottimale;
    \item Un insieme diversificato di operatori genetici e di ricerca locale;
    \item L'ottimizzazione iperparametrica con Optuna che ha prodotto la
    configurazione Tuned;
    \item Un'architettura software moderna con visualizzazione in tempo reale;
    \item Robustezza e consistenza dei risultati tra run diverse;
    \item Gap medio dello 0.84\% sul benchmark CVRPLIB, con prestazioni
    eccellenti sulle istanze di set A e B.
\end{itemize}

\section{Sviluppi Futuri}

Possibili estensioni del lavoro includono:
\begin{itemize}
    \item L'implementazione di una ricerca locale pi\`u aggressiva (Variable
    Neighborhood Search) \cite{vidal2012unified};
    \item L'uso di tecniche di diversificazione della popolazione (crowding,
    fitness sharing);
    \item L'estensione ad altre varianti del VRP (VRPTW con finestre temporali,
    VRPPD con pickup and delivery);
    \item L'implementazione di un algoritmo memetico con ricombinazione
    ottimale dei percorsi (edge assembly crossover) \cite{vidal2013hybrid};
    \item L'integrazione con un solver esatto per il post-processing delle
    route (set partitioning).
\end{itemize}

Il codice sorgente del progetto, comprensivo di backend Python, frontend React
e relazione \LaTeX, \`e disponibile nella repository del progetto.

\vspace{1cm}
\begin{center}
\rule{0.5\textwidth}{0.5mm}
\end{center}
